% gerado por scripts/build_paper_assets.py -- nao editar a mao
\begin{table}[htbp]
\centering
\small
\setlength{\tabcolsep}{4pt}
\caption{Effect of monthly minimum temperature as an exogenous covariate, under the
leakage-free climatology policy, and the labelled ceiling scenario.}
\label{tab:temperatura}
\begin{tabular}{lrrrcrr}
\toprule
Model & Without & With temp. & Gain (pp) & 95\% CI & DM cells $p<0.05$ & Ceiling \\
\midrule
SARIMA & 4.80 & 4.66 & 0.136 & $[+0.034, +0.257]$ & 1 of 6 & 4.63 \\
CatBoost & 6.59 & 6.33 & 0.258 & $[+0.073, +0.452]$ & 1 of 6 & 6.30 \\
XGBoost & 6.94 & 6.55 & 0.382 & $[+0.222, +0.543]$ & 2 of 6 & 6.55 \\
\bottomrule
\end{tabular}

\vspace{0.5em}
\begin{minipage}{\textwidth}
\footnotesize
Under the climatology policy the future covariate is the month-of-year mean recomputed
for each window from the exogenous series truncated at that window's training end, so no
value from after the training end is visible. The ceiling column replaces it with the
true observed future temperature, which leaks by construction and is reported only to
bound what a perfect weather forecast could add. All three gains have intervals
excluding zero, but none reaches the pre-declared threshold of DM significance in at
least three horizons, so the effect is reported as suggestive and not established.
Prophet and TimesFM do not accept exogenous regressors in this implementation and were
excluded from this comparison rather than being given an input they would ignore.
\end{minipage}
\end{table}
